\documentclass{report}
\usepackage{amsmath,amssymb}
\setlength{\parindent}{0mm}
\setlength{\parskip}{1em}
\begin{document}
\begin{center}
\rule{6in}{1pt} \
{\large Ray Tuminaro \\
{\bf A Quasi-Algebraic Multigrid Approach to PDE systems that include Material Interfaces and Variable Degrees-of-Freedom Per Node}}

Sandia National Laboratories \\ MS 9159 \\ PO Box 969 \\ Livermore \\ CA 94551
\\
{\tt rstumin@sandia.gov}\\
David Noble\end{center}

The application of algebraic multigrid methods to PDE systems can be
problematic, especially when the number of degrees-of-freedom per node
varies. This talk considers a multi-phase flow problem corresponding to
the movement of an air bubble within a surrounding water body. The
incompressible Navier-Stokes equations are discretized via stabilized
linear finite elements using a moving mesh that conforms to the
bubble/air interfaces. In this formulation, the number of
degrees-of-freedom at nodes on the bubble interface is greater than the
rest of the domain due to the presence of unknowns for both air and water
pressure.

To tackle this problem with multigrid, we consider an approach that
applies AMG to a scalar Laplace-like problem on the same mesh as the
original problem with the exception that edges between interface nodes
and the rest of the domain are removed. The grid transfers for this
Laplace problem are then used to build interpolation/restriction
operators for the original incompressible flow application. The overall
scheme is quasi-algebraic in that coordinates (to build the Laplace-like
problem) are required as well as a boolean table indicating which
degrees-of-freedom are defined at each nodal location. Numerical results
are presented highlighting strengths of this approach. Additionally,
alternative algorithms for the construction of grid transfers are also
described. These alternatives essentially incorporate more matrix
information and are potentially helpful when the underlying PDE
coefficients are highly variable within each isolated region.


\end{document}
